\subsection{Vertauschen von Differentiation und Integration}

% --- Beginn des Abschnitts 3.4 ---

\begin{lemma}{}
Sei $I\subset \mathbb{R}$ ein Intervall und sei $h\in C([a,b])$ (bzw. $h(x,t)$ eine Funktion, wobei $t$ im Integral als Parameter auftaucht). Dann ist die Funktion


$$
H(x)=\int_a^b h(x,t)\,dt
$$


stetig in $x$. 
% Hinweis: In der Quelle stand "t wird ansintegriert" und dies wurde als Parameter betrachtet.
\end{lemma}

\begin{beweisbox}
Der Beweis basiert darauf, dass $h(x,t)$ für jedes $t\in [a,b]$ aufgrund der Stetigkeit gleichmäßig in $x$ ist. Insbesondere gilt


$$
|H(x)-H(x_0)|\le \int_a^b |h(x,t)-h(x_0,t)|\,dt,
$$


und da $h$ gleichmäßig stetig ist (auf $ [a,b]$ beziehungsweise einem kompakten Intervall in $t$), folgt $H(x)\to H(x_0)$ für $x\to x_0$. 
% % Einige Abschätzungen oder detaillierte Beweisschritte konnten nicht vollständig rekonstruiert werden.
\end{beweisbox}

\begin{satz}{}
Sei $f\in C\bigl(I\times [a,b]\bigr)$, wobei $I\subset\mathbb{R}$ ein Intervall ist, und sei für jedes $x\in I$ die Ableitung $\frac{\partial}{\partial x}f(x,t)$ existierend und so, dass


$$
(x,t)\mapsto \frac{\partial}{\partial x}f(x,t)
$$


stetig ist auf $I\times [a,b]$. Dann definiere


$$
F(x)=\int_a^b f(x,t)\,dt.
$$


Die Funktion $F$ ist differenzierbar in $x$ mit


$$
F'(x)=\int_a^b \frac{\partial}{\partial x}f(x,t)\,dt.
$$


\end{satz}

%  Entfernt

\begin{beispielbox}
\textbf{Beispiel:} Sei


$$
F(x)=\int_0^a \sin(xt)\,dt.
$$


Dann ist $F$ differenzierbar und es gilt


$$
F'(x)=\int_0^a t\cos(xt)\,dt.
$$


% Anmerkung: Die direkte Berechnung bestätigt in konkreten Fällen, dass der Ableitungsabschluss unter dem Integral zulässig ist.
\end{beispielbox}

\begin{satz}{}
Sei $f\in C\bigl([a,b]\times[c,d]\bigr)$. Dann gilt


$$
\int_a^b\left(\int_c^d f(x,y)\,dy\right)dx = \int_c^d\left(\int_a^b f(x,y)\,dx\right)dy.
$$


\end{satz}

\begin{bemerkungbox}
Die Klammern in den obigen Integralen dienen allein der Verdeutlichung der Reihenfolge der Integration. Details zum vollständigen Satz von Fubini entnehmen Sie weiterführender Literatur.
% Einige Ausdrucksformen (z. B. Abschätzungen mit Operatornormen) konnten nicht eindeutig rekonstruiert werden.
\end{bemerkungbox}

% --- Hinweise zu weiteren, nicht eindeutig extrahierbaren Teilen ---
% Es gibt in der Quelle noch Textstellen zu Ableitungsregeln, Differenzierbarkeitsregeln (Produkt-, Ketten- und Quotientenregel) in Banachräumen sowie zu Verallgemeinerungen des Hauptsatzes.
% Diese Abschnitte enthielten auch Abschätzungen und Notationen (z. B. zur Operatornorm) und wurden hier aufgrund der unleserlichen Handschrift nur teilweise extrahiert.
% Ebenso wurden manche Details zu "Vertauschen von Limiten" und "unter dem Integral ableiten" nur in zusammengefasster Form übernommen.

% --- Ende des Abschnitts 3.4 ---
